Une corde horizontale de masse linéique \qty{\v{mu}}{g/m} est attachée à un mur à une extrémité. Elle passe sur une poulie située à \qty{\v{L}}{m} du mur, et un bloc de \qty{\v{M}}{kg} est suspendu à son autre extrémité. On produit une impulsion près de la poulie.

\begin{enumerate}[label=\alph*)]
    \item Déterminez la vitesse de l'impulsion sur la partie horizontale de la corde.
    \item L'impulsion se réfléchit sur le mur. Combien de temps s'écoule entre sa création et son retour à la poulie ?
    \item Quelle masse faudrait-il suspendre pour que l'impulsion se propage \num{\v{n}} fois plus vite ?
\end{enumerate}
